\section{Cài đặt}
\label{sec:implementation}

Nguyên mẫu được cài đặt bằng Python~3.11 dưới dạng gói \texttt{rbac-guard}.
Lớp \texttt{parser} chuẩn hóa và kiểm tra event CSV/JSON; \texttt{rbac} lưu
người dùng, role, permission trong SQLite; \texttt{rules} thực hiện ba luật
cơ sở; và \texttt{context} duy trì baseline, sinh finding và gom incident.
\texttt{scoring} tính điểm/mức độ, còn \texttt{reporting} ghi kết quả theo
kiểu thay thế tệp tạm--đổi tên để tránh một report bị ghi dở. Hàm dịch vụ
\texttt{analyze} ghép các mô-đun này thành một pipeline; CLI và giao diện web
chỉ là hai adapter gọi cùng hàm dịch vụ, nên không có hai cách diễn giải luật
khác nhau.

\subsection{Cấu hình và giao diện chạy}

CLI cài cùng gói cung cấp bốn lệnh: \texttt{init-db} khởi tạo SQLite từ seed,
\texttt{check-access} kiểm tra một permission, \texttt{analyze} chạy pipeline
và \texttt{evaluate} so alert với \texttt{expected\_label} của dataset. Tệp
\texttt{config/default.toml} là đầu vào cấu hình duy nhất cho severity, ngưỡng
password guessing, cửa sổ thời gian và trọng số ngữ cảnh; các giá trị dùng
trong thực nghiệm là $N_p=5$, $W_p=W_c=300$ giây và các ngưỡng severity
30/60/85 như Mục~\ref{sec:methodology}. Chế độ ngữ cảnh không bật mặc định;
tùy chọn \texttt{--context-risk} mới yêu cầu profiler, finding và incident.

Báo cáo cơ sở luôn gồm \texttt{alerts.csv}, \texttt{alerts.json} và
\texttt{run\_metadata.json}. Ở chế độ ngữ cảnh, chương trình còn tạo
\texttt{context\_findings.json}, \texttt{incidents.csv} và
\texttt{incidents.json}. Tệp metadata lưu số event hợp lệ và lỗi parser, nhờ
đó có thể đối chiếu report với lần chạy. Giao diện Streamlit là tùy chọn, chỉ
đọc: nó gọi cùng service, lọc alert/incident và cho tải kết quả; nó không thay
đổi RBAC hay luật phát hiện.

\subsection{Khả năng tái lập}
\label{subsec:reproducibility}

Sau khi tạo môi trường Python~3.11 và cài \texttt{pip install -e '.[dev]'},
Listing~\ref{lst:reproduce} tái sinh hai checkpoint dùng trong bài. Thư mục
tạm bảo đảm CSDL RBAC mới được seed cho mỗi lần chạy. Các đường dẫn output có
thể thay bằng một thư mục trống khác; không có JSON hay metric nào được sửa
thủ công sau khi chạy.

\begin{lstlisting}[style=cli,caption={Tái sinh kết quả baseline và context-risk.},label={lst:reproduce}]
tmpdir=$(mktemp -d)
.venv/bin/rbac-guard init-db --db "$tmpdir/rbac.db" --seed data/rbac_seed.json
.venv/bin/rbac-guard analyze --db "$tmpdir/rbac.db" --log data/logs_demo.csv \
  --config config/default.toml --output paper/risk_based_access_monitoring/artifacts/cp3-baseline
.venv/bin/rbac-guard evaluate \
  --alerts paper/risk_based_access_monitoring/artifacts/cp3-baseline/alerts.json \
  --events data/logs_demo.csv \
  --output paper/risk_based_access_monitoring/artifacts/cp3-baseline/metrics.json
.venv/bin/rbac-guard analyze --db "$tmpdir/rbac.db" --log data/logs_risk_demo.csv \
  --config config/default.toml \
  --output paper/risk_based_access_monitoring/artifacts/cp3-context --context-risk
\end{lstlisting}

Kho nguồn có kiểm thử cho parser, RBAC, luật, scoring, context, reporting,
CLI và helper web. Cổng chất lượng tái lập là
\texttt{.venv/bin/pytest -q --cov=rbac\_guard --cov-report=term-missing
--cov-fail-under=85}; coverage loại trừ hai adapter \texttt{cli.py} và
\texttt{web.py}, vốn được kiểm tra qua subprocess, helper và smoke test. Kết
quả test của checkpoint được ghi tại Mục~\ref{sec:evaluation}, còn artifact
CP3 là nguồn cho mọi số liệu trong các bảng.
